\section{Розділ~\ref{sec:pain}. Біль}

Датчики болю, два проводи, ворота в спинному мозку, низхідну ручку гучності, сенсибілізацію й захоплення уваги виміряно прямо; формули воріт і сенсибілізації --- спрощені моделі книги; правило, за яким машина вибирає гучність тривоги, --- гіпотеза.

{\footnotesize
\setlength{\tabcolsep}{3pt}\renewcommand{\arraystretch}{1.15}
\begin{longtable}{>{\raggedright}p{0.5cm}>{\raggedright}p{4.0cm}>{\raggedright}p{5.6cm}>{\raggedright}p{2.3cm}>{\raggedright\arraybackslash}p{1.7cm}}
\toprule
№ & Твердження & Дослід і що виміряно & Джерело & Статус \\
\midrule\endfirsthead
\toprule
№ & Твердження & Дослід і що виміряно & Джерело & Статус \\
\midrule\endhead
1 & Високий поріг: пороги нагрівання в різних датчиків болю від $41^\circ$ до $46$--$48^\circ$ & Запис поодиноких волокон. Шкіра мавпи: медіанний поріг повільних (C) датчиків $41^\circ$, швидких датчиків другого типу $46^\circ$, першого типу --- вище за $53^\circ$. Шкіра людини: C-датчики, що відповідають і на тиск, $40.7^\circ$, що не відповідають на тиск --- $48^\circ$. & \cite{Treede1995heat, Weidner1999cfib} & виміряно \\
2 & Датчики болю звикають набагато слабше за датчики дотику й після легкого ушкодження стають чутливішими; послаблення після сильного нагрівання виміряно в одного типу & Легкий опік шкіри ($50^\circ$, $100\,\text{с}$): через $5$--$10$~хв поріг болю в людей і поріг C-датчиків у мавпи знижено на $2$--$6^\circ$; в інших шкірних датчиків такого зсуву немає. У швидких датчиків другого типу відповідь після сильного нагрівання пригнічено. Порівняння звикання з датчиками дотику --- загальна оцінка, окремим дослідом тут не підкріплена. & \cite{LaMotte1982hyper, Treede1995heat} & виміряно \\
3 & Два проводи: швидкий $5$--$30$~м/с, повільний $0.5$--$2$~м/с; час приходу $t=L/v$~\eqref{eq:paintime} & Швидкість проведення: швидкі датчики болю мавпи в середньому $15$ і $25$~м/с (два типи); C-датчики людини $0.8$--$1.0$~м/с. При $L=1$~м це десятки мілісекунд і близько секунди. & \cite{Treede1995heat, Weidner1999cfib} & виміряно \\
4 & Біль приходить двічі: спершу швидкий і точний, потім тупий і довгий & Час реакції на нагрівання передпліччя людини: від $1100$ до $400\ms$ зі зростанням температури; сильне нагрівання волосистої шкіри дає подвійний біль, долоні --- ні. Перший біль можуть нести лише волокна, швидші за $6$~м/с, --- за латентністю йому відповідають швидкі датчики другого типу, яких у долоні немає. Розподіл ролей («де» і «наскільки погано») --- тлумачення. & \cite{CampbellLaMotte1983first, Treede1995heat} & виміряно \\
5 & Ворота: вихідний нейрон спинного мозку отримує біль і дотик, гальмівний посередник дотиком збуджується (рис.~\ref{fig:gate}) & Схему воріт запропоновано як теорію. Миші, генетична мітка й видалення груп нейронів: збуджувальні нейрони із соматостатином потрібні для механічного болю; гальмівні нейрони з динорфіном не дають входу від датчиків дотику дійти до них --- без цих нейронів легкий дотик викликає біль. & \cite{MelzackWall1965, Duan2014} & виміряно \\
6 & Дотик приглушує біль & Людина, лазерні імпульси болю на руці: одночасний дотик знижує оцінку болю й здатність розрізняти енергію імпульсів; ефект слабшає з відстанню між точкою дотику й точкою болю в межах одного шкірного сегмента. & \cite{Mancini2014touch} & виміряно \\
7 & Припущення теорії воріт: сильний біль сам вимикає гальмівного посередника, і ворота розчахуються & У розділі позначено як припущення вихідної теорії. Праця на мишах описує, як ворота не пускають дотик у канал болю; вимкнення посередника болем у ній не показано. & \cite{MelzackWall1965} & гіпотеза \\
8 & Ворота~\eqref{eq:gate}: поріг $0.18$ без дотику, $0.86$ з дотиком, $0.11$ при $g=2$; при $\nu=1$ дотик знижує вихід з $1$ до $0.25$, при $\nu=2$ не діє; дотик сам не проходить (рис.~\ref{fig:gatecurves}) & Розрахунок за \texttt{models/\allowbreak pain\_\allowbreak gate.py} з вагами з тексту відтворює всі ці числа. & \texttt{models/\allowbreak pain\_\allowbreak gate.py} & модель \\
9 & Низхідні шляхи зі стовбура змінюють підсилення воріт в обидва боки & Щур, довгастий мозок, запис під час нагрівання хвоста: одні нейрони розряджаються перед відсмикуванням («увімк»), інші замовкають («вимк»); слабка стимуляція цієї ділянки пригнічує рефлекс. Огляд: ті самі кола і полегшують, і гальмують передачу болю, опіоїди діють через них. & \cite{Fields1983rvm, Fields2004} & виміряно \\
10 & Очікування полегшення приглушує біль через опіоїдний канал & Люди з гострим болем: у тих, кому очікування полегшення зменшило біль, блокада опіоїдних рецепторів повернула біль до рівня решти; в інших блокада болю не змінювала. & \cite{Levine1978} & виміряно \\
11 & Гучність тривоги мозок вибирає за вигодою переривання & Досліду, в якому виміряно правило вибору гучності, немає. & --- & гіпотеза \\
12 & Сенсибілізація: після ушкодження вихід воріт зростає, поріг падає, почасти за рахунок спинного мозку; вона тримається й після припинення ушкоджувального стимулу & Щур: після ушкодження шкіри поріг згинального рефлексу падає, а відповідь зростає; електрофізіологічний аналіз показує, що частина зростання виникає в самому спинному мозку. Ушкоджувальний стимул викликає тривалі порушення чутливості, що переживають сам стимул. Точного часу спаду розділ не наводить. & \cite{Woolf1983} & виміряно \\
13 & Сенсибілізація~\eqref{eq:sensitization} --- додатний зворотний зв'язок; за сильної петлі тривога може защепнутися після загоєння & Випливає з моделі~\eqref{eq:sensitization} (задача наприкінці розділу). Досліду, який показав би, що стійкий біль після загоєння --- защепнута петля саме такого вигляду, немає. & виведення в розділі & гіпотеза \\
14 & Біль --- від'ємна винагорода, і навчання на ньому йде за помилкою часових різниць & Томограф, люди, ланцюжки провісників болю: активність вентрального стріатума й передньої острівцевої кори збігається із сигналами навчання за часовими різницями. Мавпа: частина \nt{DOPA}-нейронів збуджується від неприємного обдування та його провісника, інша частина гальмується. & \cite{Seymour2004, Matsumoto2009} & виміряно \\
15 & Біль перериває поточну роботу & Люди, електричний больовий стимул і звукові проби: відповідь на пробу через $100\ms$ після початку болю помітно сповільнена, через $1.5\,\text{с}$ різниці з контролем уже немає. Огляд зводить такі досліди в модель захоплення уваги. & \cite{Crombez1996disrupt, EcclestonCrombez1999} & виміряно \\
16 & Відсмикування замикається в спинному мозку & Щур: нейрони глибоких шарів заднього рогу повторюють просторовий візерунок рефлексу відсмикування окремих м'язів; їх не вдається збудити антидромно з шийного відділу, тобто це спінальні посередники. & \cite{Schouenborg1995wr} & виміряно \\
\bottomrule
\end{longtable}}
